\section{序論}
今回の観測窓では、GoogleのGemini 4 Neonに関する継続的な検証や利用感想が目立った。モデルの性能や価格に加え、コーディングエージェントでの実用性を巡る議論が活発である。AnthropicのIPOを巡る報道や、Runwayによるロボティクス向けモデルの公開も関心を集めた。

研究面では、Transformerの内部状態を拡張する言語モデル、AIを研究開発へ組み込む動き、タンパク質構造予測に関する話題が観測された。オープンモデルの小型化やエージェント用途への最適化も継続的なテーマとなっている。

本号はGrokの収集結果8件を記載順に記事化したものである。Grokの網羅性評価はPartialであり、窓内のすべての議論を代表するものではない。